\subsection{ELIMINATION THEOREM}

PROPOSITION 9. Let \(\mathbf{S}_1\) and \(\mathbf{S}_2\) be two disjoint sets and \(d\) a mapping of \(\mathbf{S}_1 \times \mathbf{S}_2\) into \(\mathrm{L}(\mathbf{S}_2)\) . Let \(\mathfrak{g}\) be the quotient Lie algebra of \(\mathrm{L}(\mathbf{S}_1 \cup \mathbf{S}_2)\) by the ideal generated by the elements \([s_1, s_2] - d(s_1, s_2)\) with \(s_1 \in \mathbf{S}_1\) , \(s_2 \in \mathbf{S}_2\) ; let \(\psi\) be the canonical mapping of \(\mathrm{L}(\mathbf{S}_1 \cup \mathbf{S}_2)\) onto \(\mathfrak{g}\) .

\begin{enumerate}
\def\labelenumi{(\alph{enumi})}
\item
  For \(i = 1,2\) , the restriction \(\phi_{i}\) of \(\psi\) to \(S_{i}\) can be extended to an isomorphism of \(\mathbf{L}(S_i)\) onto a subalgebra \(\mathfrak{a}_i\) of \(\mathfrak{g}\) .
\item
  \(\mathfrak{g} = \mathfrak{a}_1 + \mathfrak{a}_2, \mathfrak{a}_1 \cap \mathfrak{a}_2 = \{0\}\) and \(\mathfrak{a}_2\) is an ideal of \(\mathfrak{g}\) .
\end{enumerate}

For \(i = 1,2\) , let \(\psi_{i}\) denote the homomorphism of \(\mathbf{L}(\mathbf{S}_i)\) into \(\mathfrak{g}\) which extends \(\phi_i\) and \(\mathfrak{a}_i\) its image. Clearly \(\phi_i(\mathbf{S}_i)\) generates \(\mathfrak{a}_i\) .

Let \(s_1 \in S_1\) ; we write \(D = \operatorname{ad} \phi_1(s_1)\) . The derivation \(D\) of \(g\) maps \(\phi_2(S_2)\) into \(a_2\) by the relation

\[
[ \phi_ {1} (s _ {1}), \phi_ {2} (s _ {2}) ] = \psi_ {2} (d (s _ {1}, s _ {2})) \quad \text { for } s _ {2} \in S _ {2};
\]

as the subalgebra \(\mathfrak{a}_2\) of \(\mathfrak{g}\) is generated by \(\phi_2(S_2)\) , therefore \(D(\mathfrak{a}_2) \subset \mathfrak{a}_2\) . The set of \(x \in \mathfrak{g}\) such that ad \(x\) leaves \(\mathfrak{a}_2\) invariant is a Lie subalgebra of \(\mathfrak{g}\) which contains \(\phi_1(S_1)\) by the above and hence also \(\mathfrak{a}_1\) . Hence

\[
\left[ \mathfrak {a} _ {1}, \mathfrak {a} _ {2} \right] \subset \mathfrak {a} _ {2}.\tag{23}
\]

Therefore \(a_1 + a_2\) is a Lie subalgebra of \(g\) and, as it contains the generating set \(\phi_1(S_1) \cup \phi_2(S_2)\) ,

\[
a _ {1} + a _ {2} = g.\tag{24}
\]

For all \(s_1 \in S_1\) there exists a derivation \(D_{s_1}\) of \(L(S_2)\) such that

\[
\mathrm{D} _ {s _ {1}} (s _ {2}) = d (s _ {1}, s _ {2})
\]

for all \(s_2\) in \(S_2\) (no. 8, Corollary to Proposition 8). The mapping \(s_1 \mapsto D_{s_1}\) can be extended to a homomorphism \(D\) of \(L(S_1)\) into the Lie algebra of derivations of \(L(S_2)\) . Let \(\mathfrak{h}\) be the semi-direct product of \(L(S_1)\) by \(L(S_2)\) corresponding to \(D\) (Chapter I, § 1, no. 8). As a module \(\mathfrak{h}\) is equal to \(L(S_1) \times L(S_2)\) and in particular

\[
[ (s _ {1}, 0), (0, s _ {2}) ] = (0, d (s _ {1}, s _ {2}))\tag{25}
\]

for \(s_{1} \in S_{1}\) and \(s_{2} \in S_{2}\) .

From (25) we deduce the existence of a homomorphism \(f\) of \(\mathfrak{g}\) into \(\mathfrak{h}\) such that \(f(\phi_1(s_1)) = (s_1, 0)\) and \(f(\phi_2(s_2)) = (0, s_2)\) for \(s_1 \in S_1\) and \(s_2 \in S_2\) . We deduce immediately the relation

\[
f (\psi_ {1} (a _ {1}) + \psi_ {2} (a _ {2})) = (a _ {1}, a _ {2})\tag{26}
\]

for \(a_1 \in \mathrm{L}(\mathrm{S}_1)\) and \(a_2 \in \mathrm{L}(\mathrm{S}_2)\) .

Relation (26) shows that \(\psi_{1}\) and \(\psi_{2}\) are injective and that \(a_{1} \cap a_{2} = \{0\}\) . Formulae (23) and (24) then imply the proposition.

PROPOSITION 10 (elimination theorem). Let X be a set, S a subset of X and T the set of sequences \((s_1, \ldots, s_n, x)\) with \(n \geqslant 0\) , \(s_1, \ldots, s_n\) in S and \(x\) in \(\mathbf{X} - \mathbf{S}\) .

\begin{enumerate}
\def\labelenumi{(\alph{enumi})}
\item
  The module \(\mathbf{L}(\mathbf{X})\) is the direct sum of the subalgebra \(\mathbf{L}(\mathbf{S})\) of \(\mathbf{L}(\mathbf{X})\) and the ideal \(\mathfrak{a}\) of \(\mathbf{L}(\mathbf{X})\) generated by \(\mathbf{X} - \mathbf{S}\) .
\item
  There exists a Lie algebra isomorphism \(\phi\) of L(T) onto \(\mathfrak{a}\) which maps \((s_1, \ldots, s_n, x)\) to (ad \(s_1 \circ \cdots \circ\) ad \(s_n\) ) ( \(x\) ).
\end{enumerate}

Let \(g\) be the Lie algebra constructed as in Proposition 9 given

\[
\mathrm{S} _ {1} = \mathrm{S}, \quad \mathrm{S} _ {2} = \mathrm{T}, \quad d (s, t) = (s, s _ {1}, \dots , s _ {n}, x) \in \mathrm{T} \subset \mathrm{L} (\mathrm{T})
\]

for \(t = (s_1, \ldots, s_n, x)\) in \(T\) and \(s \in S_1\) . We identify \(L(S)\) and \(L(T)\) with their canonical images in \(g\) (Proposition 9 (a)).

Let \(\psi\) be the mapping \((s_1, \ldots, s_n, x) \mapsto (\operatorname{ad} s_1 \circ \cdots \circ \operatorname{ad} s_n)(x)\) of T into \(\mathrm{L}(\mathbf{X})\) . Obviously \(\psi(d(s, t)) = [s, \psi(t)]\) for \(s \in S\) and \(t \in T\) and hence there exists a homomorphism \(\alpha: g \to L(X)\) whose restriction to \(S\) is the identity and whose restriction to \(T\) is \(\psi\) . Now \(X - S \subset T\) , whence there is a homomorphism \(\beta: L(X) \to g\) whose restriction to \(X = S \cup (X - S)\) is the identity.

We show that \(\alpha\) is an isomorphism and \(\beta\) the inverse isomorphism. As \(\psi(x) = x\) for \(x\) in \(\mathbf{X} - \mathbf{S}\) , we see that \(\alpha \circ \beta\) coincides with the identity on \(\mathbf{X}\) , whence \(\alpha \circ \beta = \mathrm{Id}_{\mathrm{L}(\mathbf{X})}\) . On the other hand, \([s, t] = d(s, t)\) in \(\mathfrak{g}\) for \(s \in \mathbb{S}\) , \(t \in \mathbb{T}\) by construction; it follows that \(t = (s_1, \ldots, s_n, x)\) is equal in \(\mathfrak{g}\) to \((\operatorname{ad} s_1 \circ \cdots \circ \operatorname{ad} s_n)(x)\) whence \(t = \beta(\alpha(t))\) . As \(\beta(\alpha(s)) = s\) for \(s \in \mathbb{S}\) and \(\mathbb{S} \cup \mathbb{T}\) generates \(\mathfrak{g}, \beta \circ \alpha = \mathrm{Id}_{\mathfrak{g}}\) .

As \(\alpha\) is an isomorphism of \(g\) onto \(L(X)\) , Proposition 9 shows that the restriction of \(\alpha\) to \(\mathrm{L(T)}\) is an isomorphism \(\phi\) of \(\mathrm{L(T)}\) onto an ideal \(\mathfrak{b}\) of \(\mathrm{L(X)}\) such that the module \(\mathrm{L(X)}\) is the direct sum of \(\mathrm{L(S)}\) and \(\mathfrak{b}\) . Obviously

\[
\phi (s _ {1}, \dots , s _ {n}, x) = (\operatorname{ad} s _ {1} \circ \dots \circ \operatorname{ad} s _ {n}) (x)
\]

for \((s_1, \ldots, s_n, x)\) in T.

Hence \(\phi(T) \subset \mathfrak{a}\) , whereas \(\mathfrak{b} \subset \mathfrak{a}\) since \(\phi(T)\) generates the subalgebra \(\mathfrak{b}\) of \(L(X)\) . But \(\mathfrak{b}\) is an ideal and \(X - S \subset \phi(T) \subset \mathfrak{b}\) , whence \(\mathfrak{a} \subset \mathfrak{b}\) .

COROLLARY. Let \(y \in \mathbf{X}\) . The free Lie algebra \(\mathbf{L}(\mathbf{X})\) is the direct sum of the free submodule \(\mathbf{K}.y\) and the Lie subalgebra admitting as basic family the family of \(((\operatorname{ad}y)^n.z)\) where \(n \geqslant 0\) and \(z \in \mathbf{X} - \{y\}\) .

It suffices to put \(S = \{y\}\) in Proposition 10.

